\documentclass[11pt,a4paper]{article}

\usepackage[T1]{fontenc}
\usepackage[utf8]{inputenc}
\usepackage{lmodern}
\usepackage{microtype}
\usepackage{amsmath}
\usepackage[a4paper, margin=0.85in, top=0.9in, bottom=0.9in]{geometry}
\usepackage{xcolor}
\usepackage{booktabs}
\usepackage{enumitem}
\usepackage{fancyhdr}
\usepackage{titlesec}
\usepackage{listings}
\usepackage[colorlinks=true, linkcolor=black, urlcolor=black]{hyperref}

% --- Academic Monochrome-Slate Palette ---
\definecolor{ink}{HTML}{111827}
\definecolor{muted}{HTML}{4B5563}
\definecolor{rulegray}{HTML}{D1D5DB}
\definecolor{codebg}{HTML}{F9FAFB}
\definecolor{kwblue}{HTML}{1E3A5F}
\definecolor{strgray}{HTML}{374151}
\definecolor{cmtgray}{HTML}{6B7280}

% --- Typography & Spacing ---
\setlength{\parindent}{0pt}
\setlength{\parskip}{0.45\baselineskip}
\setlist[itemize]{nosep, leftmargin=1.35em, itemsep=2pt, topsep=3pt}

% --- Section Styling ---
\titleformat{\section}
  {\large\bfseries\color{ink}}
  {}
  {0em}
  {}
  [\vspace{2pt}{\color{rulegray}\titlerule[0.5pt]}]
\titlespacing*{\section}{0pt}{1.4\baselineskip}{0.5\baselineskip}

% --- Running Headers & Footers ---
\pagestyle{fancy}
\fancyhf{}
\fancyhead[L]{\footnotesize\scshape\color{muted} Operating System Laboratory Copy}
\fancyhead[R]{\footnotesize\color{muted} Ronit Choudhury \textbar\ IS-241-844-0008}
\fancyfoot[C]{\footnotesize\color{muted}\thepage}
\setlength{\headheight}{14pt}
\renewcommand{\headrulewidth}{0.4pt}
\renewcommand{\footrulewidth}{0pt}

% --- Source Code Listing Styling ---
\lstdefinestyle{academicC}{
    language=C,
    basicstyle=\ttfamily\footnotesize\color{ink},
    keywordstyle=\bfseries\color{kwblue},
    commentstyle=\itshape\color{cmtgray},
    stringstyle=\color{strgray},
    numberstyle=\tiny\color{cmtgray},
    numbers=left,
    numbersep=9pt,
    stepnumber=1,
    frame=lines,
    framerule=0.45pt,
    rulecolor=\color{rulegray},
    backgroundcolor=\color{codebg},
    xleftmargin=1.6em,
    framexleftmargin=1.4em,
    aboveskip=0.6\baselineskip,
    belowskip=0.8\baselineskip,
    tabsize=4,
    showstringspaces=false,
    breaklines=true,
    breakatwhitespace=true,
    keepspaces=true,
    columns=flexible,
    captionpos=t
}
\lstset{style=academicC}

% --- Compact Experiment Metadata Helper ---
\newcommand{\exptheader}[3]{%
    \section[#1: #2]{#1 --- #2}%
    \label{#3}%
}

\begin{document}

% ============================================================================
% TITLE PAGE
% ============================================================================
\begin{titlepage}
    \thispagestyle{empty}
    \centering
    \vspace*{1.6cm}

    {\small\scshape\color{muted} Gauhati University \par}
    \vspace{0.25cm}
    {\normalsize\scshape\color{ink} Department of IT \par}

    \vspace{1.2cm}
    {\color{rulegray}\rule{0.45\textwidth}{0.6pt}\par}
    \vspace{1.2cm}

    {\LARGE\bfseries\color{ink} Operating System Laboratory Copy \par}
    \vspace{0.45cm}
    {\normalsize\itshape\color{muted}
        System Calls, Process Control, CPU Scheduling, Synchronization, and Deadlock Avoidance
    \par}

    \vspace{1.2cm}
    {\color{rulegray}\rule{0.45\textwidth}{0.6pt}\par}

    \vfill

    \begin{tabular}{@{}ll@{}}
        \toprule
        \textbf{Name}        & Ronit Choudhury \\
        \textbf{Roll No.}    & IS-241-844-0008 \\
        \textbf{Department}  & Department of IT \\
        \textbf{Institution} & Gauhati University \\
        \textbf{Language}    & C (ISO C / POSIX API) \\
        \bottomrule
    \end{tabular}

    \vspace{1.8cm}
\end{titlepage}

% ============================================================================
% INDEX PAGE
% ============================================================================
\newpage
\thispagestyle{plain}
\section*{Index}
\vspace{0.2cm}

\begin{center}
\small
\begin{tabular}{@{}lllr@{}}
    \toprule
    \textbf{Expt.\ No.} & \textbf{Experiment Title} & \textbf{Source File} & \textbf{Page} \\
    \midrule
    2.1  & \hyperref[exp:2.1]{Directory Listing via Interactive Input}          & \texttt{expt\_2\_1.c} & \pageref{exp:2.1} \\
    2.2  & \hyperref[exp:2.2]{Process Creation and Process Identification}      & \texttt{expt\_2\_2.c} & \pageref{exp:2.2} \\
    3.1  & \hyperref[exp:3.1]{Simulation of the UNIX \texttt{ls} Command}       & \texttt{expt\_3\_1.c} & \pageref{exp:3.1} \\
    3.2  & \hyperref[exp:3.2]{Simulation of the UNIX \texttt{grep} Command}     & \texttt{expt\_3\_2.c} & \pageref{exp:3.2} \\
    3.3  & \hyperref[exp:3.3]{Pattern Matching and Occurrence Counting}           & \texttt{expt\_3\_3.c} & \pageref{exp:3.3} \\
    5(a) & \hyperref[exp:5.a]{Non-Preemptive Priority CPU Scheduling}           & \texttt{expt\_5\_a.c} & \pageref{exp:5.a} \\
    5(b) & \hyperref[exp:5.b]{Round Robin (RR) CPU Scheduling}                  & \texttt{expt\_5\_b.c} & \pageref{exp:5.b} \\
    5(c) & \hyperref[exp:5.c]{First-Come, First-Served (FCFS) CPU Scheduling}     & \texttt{expt\_5\_c.c} & \pageref{exp:5.c} \\
    5(d) & \hyperref[exp:5.d]{Shortest Job First (SJF) CPU Scheduling}          & \texttt{expt\_5\_d.c} & \pageref{exp:5.d} \\
    6    & \hyperref[exp:6]{Producer--Consumer Problem Using Semaphores}        & \texttt{expt\_6.c}    & \pageref{exp:6} \\
    8    & \hyperref[exp:8]{Deadlock Avoidance Using Banker's Algorithm}          & \texttt{expt\_8.c}    & \pageref{exp:8} \\
    \bottomrule
\end{tabular}
\end{center}

\newpage

% ============================================================================
% EXPERIMENT 2.1
% ============================================================================
\exptheader{Experiment 2.1}{Directory Listing via Interactive Input}{exp:2.1}

This program prompts the user for a filesystem path and enumerates all entries (files, subdirectories, \texttt{.} and \texttt{..}) residing in that directory. It opens a POSIX directory stream via \texttt{opendir()}, iterates sequentially through \texttt{struct dirent} records using \texttt{readdir()}, and releases the underlying descriptor with \texttt{closedir()}.

\begin{itemize}
    \item \textbf{POSIX Directory Streams:} Abstraction over kernel directory entries (\texttt{DIR *} and \texttt{struct dirent}).
    \item \textbf{System-Level Traversal:} Relies on underlying directory-read system calls (e.g., \texttt{getdents} on Linux).
    \item \textbf{Error Reporting:} Uses \texttt{perror()} and \texttt{stderr} to diagnose invalid or inaccessible paths.
\end{itemize}

\begin{lstlisting}[title={\texttt{expt\_2\_1.c}}]
#include <stdio.h>
#include <dirent.h>

int main(void)
{
    char path[256];
    struct dirent *entry;
    DIR *dir;

    /* Read up to 255 characters to prevent buffer overflow on path input. */
    printf("Enter directory name: ");
    if (scanf("%255s", path) != 1) {
        fprintf(stderr, "No directory name given.\n");
        return 1;
    }

    /* Open the directory stream; returns NULL if path is invalid/inaccessible. */
    dir = opendir(path);
    if (dir == NULL) {
        perror(path);
        return 1;
    }

    /* Sequentially fetch each directory entry until end-of-stream (NULL). */
    while ((entry = readdir(dir)) != NULL)
        printf("%s\n", entry->d_name);

    /* Release the directory stream and its underlying file descriptor. */
    closedir(dir);
    return 0;
}
\end{lstlisting}

% ============================================================================
% EXPERIMENT 2.2
% ============================================================================
\exptheader{Experiment 2.2}{Process Creation and Process Identification}{exp:2.2}

This program demonstrates OS process spawning and process identification across POSIX and Windows environments. On POSIX systems, \texttt{fork()} duplicates the calling process and inspects the return value (\texttt{0} in the child, child PID in the parent); on Windows, \texttt{\_spawnv()} creates a new child instance while both processes query their Process IDs via \texttt{getpid()}.

\begin{itemize}
    \item \textbf{Process Creation Models:} POSIX process duplication (\texttt{fork}) versus Windows process spawning (\texttt{\_spawnv}).
    \item \textbf{Process Control Block (PCB):} Each process receives a distinct kernel-assigned Process Identifier (PID).
    \item \textbf{I/O Buffer Consistency:} \texttt{fflush(stdout)} prevents buffered parent output from being duplicated across \texttt{fork()}.
\end{itemize}

\begin{lstlisting}[title={\texttt{expt\_2\_2.c}}]
#include <stdio.h>
#include <stdlib.h>
#include <string.h>

#ifdef _WIN32
#include <process.h>
#define getpid _getpid
#else
#include <unistd.h>
#endif

int main(int argc, char *argv[])
{
    /* Windows starts a fresh copy of this program as the child. */
    if (argc > 1 && strcmp(argv[1], "child") == 0) {
        printf("Child process ID: %d\n", getpid());
        return 0;
    }

    /* Flush stdout before creating the child to avoid duplicated buffered output. */
    printf("Parent process ID: %d\n", getpid());
    fflush(stdout);

#ifdef _WIN32
    const char *child_args[] = { argv[0], "child", NULL };

    /* Synchronous process creation on Windows (_P_WAIT waits for completion). */
    if (_spawnv(_P_WAIT, argv[0], child_args) == -1) {
        perror("spawn");
        return 1;
    }
#else
    {
        /* Duplicate the parent process; fork() returns twice on success. */
        pid_t pid = fork();
        if (pid < 0) {
            perror("fork");
            return 1;
        }
        /* Execution branch unique to the newly created child process. */
        if (pid == 0)
            printf("Child process ID: %d\n", getpid());
    }
#endif

    return 0;
}
\end{lstlisting}

% ============================================================================
% EXPERIMENT 3.1
% ============================================================================
\exptheader{Experiment 3.1}{Simulation of the UNIX \texttt{ls} Command}{exp:3.1}

This utility replicates the core listing behaviour of the UNIX \texttt{ls} command by accepting a target directory path through command-line arguments. After verifying \texttt{argc == 2}, it opens the directory stream with \texttt{opendir()}, prints each entry's \texttt{d\_name} field via \texttt{readdir()}, and closes the stream cleanly.

\begin{itemize}
    \item \textbf{Command-Line Interface:} Argument vector (\texttt{argc}, \texttt{argv}) parsing standard to UNIX utilities.
    \item \textbf{Directory Stream API:} Sequential traversal using \texttt{opendir()}, \texttt{readdir()}, and \texttt{closedir()}.
    \item \textbf{Namespace \& Metadata:} Accessing filename strings stored within \texttt{struct dirent}.
\end{itemize}

\begin{lstlisting}[title={\texttt{expt\_3\_1.c}}]
#include <stdio.h>
#include <dirent.h>

int main(int argc, char *argv[])
{
    struct dirent *entry;
    DIR *dir;

    /* Enforce exact argument count: program name plus target directory. */
    if (argc != 2) {
        fprintf(stderr, "Usage: %s <directory>\n", argv[0]);
        return 1;
    }

    /* Open directory stream and report OS-level errors (e.g., ENOENT) via perror. */
    dir = opendir(argv[1]);
    if (dir == NULL) {
        perror(argv[1]);
        return 1;
    }

    /* Traverse all directory entries until readdir() reaches end-of-directory. */
    printf("Contents of %s:\n", argv[1]);
    while ((entry = readdir(dir)) != NULL)
        printf("%s\n", entry->d_name);

    /* Close the stream to free kernel directory handle resources. */
    closedir(dir);
    return 0;
}
\end{lstlisting}

% ============================================================================
% EXPERIMENT 3.2
% ============================================================================
\exptheader{Experiment 3.2}{Simulation of the UNIX \texttt{grep} Command}{exp:3.2}

This program simulates the line-filtering behaviour of UNIX \texttt{grep} by scanning a text file for a specified search word. It opens the file via \texttt{fopen()}, reads bounded lines with \texttt{fgets()}, removes trailing newline characters using \texttt{strcspn()}, and prints every line containing the substring returned by \texttt{strstr()} alongside its filename and line number.

\begin{itemize}
    \item \textbf{Buffered File I/O:} Stream-oriented reading with \texttt{fopen()}, \texttt{fgets()}, and \texttt{fclose()}.
    \item \textbf{Line-Oriented Filter Model:} Sequential text processing characteristic of UNIX filter tools.
    \item \textbf{Bounded Input Handling:} Fixed-size line buffer (\texttt{MAX\_LINE}) prevents stack overflow on long input.
\end{itemize}

\begin{lstlisting}[title={\texttt{expt\_3\_2.c}}]
#include <stdio.h>
#include <string.h>

#define MAX_LINE 1024

int main(int argc, char *argv[])
{
    char line[MAX_LINE];
    FILE *fp;
    int line_no = 0;

    /* Expect two CLI arguments: <file> and <word>. */
    if (argc != 3) {
        fprintf(stderr, "Usage: %s <file> <word>\n", argv[0]);
        return 1;
    }

    /* Open the target file in read-only text mode. */
    fp = fopen(argv[1], "r");
    if (fp == NULL) {
        perror(argv[1]);
        return 1;
    }

    /* Read the file sequentially line-by-line up to MAX_LINE - 1 bytes. */
    while (fgets(line, sizeof line, fp) != NULL) {
        line_no++;
        line[strcspn(line, "\n")] = '\0';   /* remove trailing newline */

        /* Print formatted match if the search word occurs anywhere in the line. */
        if (strstr(line, argv[2]) != NULL)
            printf("%s:%d: %s\n", argv[1], line_no, line);
    }

    fclose(fp);
    return 0;
}
\end{lstlisting}

% ============================================================================
% EXPERIMENT 3.3
% ============================================================================
\exptheader{Experiment 3.3}{Pattern Matching and Occurrence Counting}{exp:3.3}

Extending the \texttt{grep} simulation, this program locates all lines matching a non-empty pattern and computes the total number of non-overlapping occurrences across the file. For each line read via \texttt{fgets()}, an inner loop repeatedly invokes \texttt{strstr()} and advances the scan pointer by \texttt{strlen(pattern)} bytes.

\begin{itemize}
    \item \textbf{Sequential File Scanning:} Single-pass buffered traversal over text streams.
    \item \textbf{Non-Overlapping Substring Counting:} Pointer arithmetic advances past each matched token.
    \item \textbf{Defensive Input Validation:} Rejects empty patterns to avoid non-terminating scan loops.
\end{itemize}

\begin{lstlisting}[title={\texttt{expt\_3\_3.c}}]
#include <stdio.h>
#include <string.h>

#define MAX_LINE 1024

int main(int argc, char *argv[])
{
    char line[MAX_LINE];
    const char *pattern;
    FILE *fp;
    int line_no = 0;
    int occurrences = 0;

    if (argc != 3) {
        fprintf(stderr, "Usage: %s <file> <pattern>\n", argv[0]);
        return 1;
    }

    /* Reject an empty pattern to prevent an infinite zero-advance loop. */
    pattern = argv[2];
    if (*pattern == '\0') {
        fprintf(stderr, "Pattern cannot be empty.\n");
        return 1;
    }

    fp = fopen(argv[1], "r");
    if (fp == NULL) {
        perror(argv[1]);
        return 1;
    }

    /* Process each line and find all non-overlapping occurrences within it. */
    while (fgets(line, sizeof line, fp) != NULL) {
        char *match = line;
        int line_occurrences = 0;

        line_no++;
        while ((match = strstr(match, pattern)) != NULL) {
            line_occurrences++;
            occurrences++;
            match += strlen(pattern);   /* advance past matched pattern */
        }

        /* Print the line only once even if the pattern appeared multiple times. */
        if (line_occurrences > 0) {
            line[strcspn(line, "\n")] = '\0';
            printf("%s:%d: %s\n", argv[1], line_no, line);
        }
    }

    fclose(fp);
    printf("Total occurrences of \"%s\": %d\n", pattern, occurrences);

    return 0;
}
\end{lstlisting}

% ============================================================================
% EXPERIMENT 5(a)
% ============================================================================
\exptheader{Experiment 5(a)}{Non-Preemptive Priority CPU Scheduling}{exp:5.a}

This program simulates non-preemptive priority scheduling for $n$ processes arriving simultaneously at $t = 0$, where smaller integer values represent higher priority. Processes are sorted via \texttt{qsort()} using a deterministic comparator (\texttt{by\_priority}) that resolves priority ties by PID, after which waiting and turnaround times are computed sequentially.

\begin{itemize}
    \item \textbf{Priority Scheduling:} CPU allocation based on static integer priority rankings.
    \item \textbf{Performance Metrics:} Computes Waiting Time ($\text{WT}$) and Turnaround Time ($\text{TAT} = \text{WT} + \text{BT}$).
    \item \textbf{Starvation \& Aging:} Low-priority jobs risk indefinite blocking unless aging gradually elevates priority.
\end{itemize}

\begin{lstlisting}[title={\texttt{expt\_5\_a.c}}]
#include <stdio.h>
#include <stdlib.h>

#define MAX_PROCESSES 10

/* Process Control Block (PCB) representation for scheduling simulation. */
struct process {
    int pid;
    int priority;
    int burst;
    int wait;
    int turnaround;
};

/* qsort comparator: orders by priority ascending, then PID ascending. */
static int by_priority(const void *a, const void *b)
{
    const struct process *x = a;
    const struct process *y = b;

    if (x->priority != y->priority)
        return x->priority - y->priority;   /* smaller number = higher priority */
    return x->pid - y->pid;                 /* keep input order for ties */
}

int main(void)
{
    struct process p[MAX_PROCESSES];
    int n, i, elapsed = 0;
    int total_wait = 0;
    int total_turnaround = 0;

    printf("Priority scheduling\nEnter number of processes: ");
    if (scanf("%d", &n) != 1 || n < 1 || n > MAX_PROCESSES) {
        fprintf(stderr, "Enter a number between 1 and %d.\n", MAX_PROCESSES);
        return 1;
    }

    for (i = 0; i < n; i++) {
        p[i].pid = i + 1;
        printf("Process %d (burst time priority): ", p[i].pid);
        if (scanf("%d %d", &p[i].burst, &p[i].priority) != 2 ||
            p[i].burst < 1) {
            fprintf(stderr, "Invalid input.\n");
            return 1;
        }
    }

    /* Order processes from highest priority (lowest integer) to lowest. */
    qsort(p, n, sizeof p[0], by_priority);

    /* Non-preemptive dispatch: each process runs to completion upon selection. */
    for (i = 0; i < n; i++) {
        p[i].wait = elapsed;                  /* WT = start time (arrival = 0) */
        p[i].turnaround = elapsed + p[i].burst; /* TAT = WT + Burst Time */
        elapsed += p[i].burst;
        total_wait += p[i].wait;
        total_turnaround += p[i].turnaround;
    }

    printf("\nPID  Priority  Burst  Wait  Turnaround\n");
    for (i = 0; i < n; i++)
        printf("%3d  %8d  %5d  %4d  %10d\n",
               p[i].pid, p[i].priority, p[i].burst,
               p[i].wait, p[i].turnaround);

    printf("\nTotal waiting time    : %d\n", total_wait);
    printf("Average waiting time  : %.2f\n", (double)total_wait / n);
    printf("Total turnaround time : %d\n", total_turnaround);
    printf("Average turnaround    : %.2f\n", (double)total_turnaround / n);

    return 0;
}
\end{lstlisting}

% ============================================================================
% EXPERIMENT 5(b)
% ============================================================================
\exptheader{Experiment 5(b)}{Round Robin (RR) CPU Scheduling}{exp:5.b}

This program simulates preemptive Round Robin scheduling across $n$ processes using a user-supplied time quantum. It cycles through unfinished processes, granting each at most one quantum (\texttt{min(remaining, quantum)}), emitting a chronological Gantt chart trace, and recording turnaround and waiting times upon completion.

\begin{itemize}
    \item \textbf{Time-Sharing Preemption:} Allocates fixed CPU time slices in cyclic FIFO order.
    \item \textbf{Quantum Trade-off:} Small quanta increase context-switch overhead; large quanta degenerate to FCFS.
    \item \textbf{Fairness Guarantee:} Bounded response time ensures starvation-free execution across all ready tasks.
\end{itemize}

\begin{lstlisting}[title={\texttt{expt\_5\_b.c}}]
#include <stdio.h>

#define MAX_PROCESSES 10

/* Process descriptor tracking total burst and remaining CPU time. */
struct process {
    int pid;
    int burst;
    int remaining;
    int wait;
    int turnaround;
};

int main(void)
{
    struct process p[MAX_PROCESSES];
    int n, quantum, i;
    int time = 0, completed = 0;
    int total_wait = 0;
    int total_turnaround = 0;

    printf("Round Robin scheduling\nEnter number of processes: ");
    if (scanf("%d", &n) != 1 || n < 1 || n > MAX_PROCESSES) {
        fprintf(stderr, "Enter a number between 1 and %d.\n", MAX_PROCESSES);
        return 1;
    }

    printf("Enter time quantum: ");
    if (scanf("%d", &quantum) != 1 || quantum < 1) {
        fprintf(stderr, "Time quantum must be a positive number.\n");
        return 1;
    }

    for (i = 0; i < n; i++) {
        p[i].pid = i + 1;
        printf("Burst time of process %d: ", p[i].pid);
        if (scanf("%d", &p[i].burst) != 1 || p[i].burst < 1) {
            fprintf(stderr, "Burst time must be a positive number.\n");
            return 1;
        }
        p[i].remaining = p[i].burst;
        p[i].wait = 0;
        p[i].turnaround = 0;
    }

    /* Circular ready-queue simulation until all n processes finish. */
    printf("\nGantt chart:\n");
    while (completed < n) {
        for (i = 0; i < n; i++) {
            int slice;

            if (p[i].remaining == 0)
                continue;

            /* Grant at most one quantum (or remaining burst if smaller). */
            slice = (p[i].remaining < quantum) ? p[i].remaining : quantum;
            printf("P%d [%d - %d]\n", p[i].pid, time, time + slice);

            time += slice;
            p[i].remaining -= slice;

            /* Record completion metrics once the process finishes its burst. */
            if (p[i].remaining == 0) {
                p[i].turnaround = time;                 /* arrival = 0 */
                p[i].wait = time - p[i].burst;          /* WT = TAT - BT */
                total_wait += p[i].wait;
                total_turnaround += p[i].turnaround;
                completed++;
            }
        }
    }

    printf("\nPID  Burst  Wait  Turnaround\n");
    for (i = 0; i < n; i++)
        printf("%3d  %5d  %4d  %10d\n",
               p[i].pid, p[i].burst, p[i].wait, p[i].turnaround);

    printf("\nTotal waiting time    : %d\n", total_wait);
    printf("Average waiting time  : %.2f\n", (double)total_wait / n);
    printf("Total turnaround time : %d\n", total_turnaround);
    printf("Average turnaround    : %.2f\n", (double)total_turnaround / n);

    return 0;
}
\end{lstlisting}

% ============================================================================
% EXPERIMENT 5(c)
% ============================================================================
\exptheader{Experiment 5(c)}{First-Come, First-Served (FCFS) CPU Scheduling}{exp:5.c}

This program implements non-preemptive First-Come, First-Served scheduling, dispatching processes strictly in their arrival order. Each process's waiting time equals the cumulative burst times of all preceding processes ($\text{WT}_i = \text{WT}_{i-1} + \text{BT}_{i-1}$), and turnaround time equals $\text{WT}_i + \text{BT}_i$.

\begin{itemize}
    \item \textbf{FIFO Ready Queue:} Simplest non-preemptive scheduling discipline with minimal scheduling overhead.
    \item \textbf{Convoy Effect:} Short CPU bursts suffer high waiting times when queued behind a long CPU-bound job.
    \item \textbf{Starvation-Free Execution:} Every queued process is guaranteed eventual CPU dispatch.
\end{itemize}

\begin{lstlisting}[title={\texttt{expt\_5\_c.c}}]
#include <stdio.h>

#define MAX_PROCESSES 10

struct process {
    int pid;
    int burst;
    int wait;
    int turnaround;
};

int main(void)
{
    struct process p[MAX_PROCESSES];
    int n, i;
    int total_wait = 0;
    int total_turnaround = 0;

    printf("FCFS scheduling\nEnter number of processes: ");
    if (scanf("%d", &n) != 1 || n < 1 || n > MAX_PROCESSES) {
        fprintf(stderr, "Enter a number between 1 and %d.\n", MAX_PROCESSES);
        return 1;
    }

    /* Compute WT and TAT inline in arrival order as a prefix sum of bursts. */
    for (i = 0; i < n; i++) {
        p[i].pid = i + 1;
        printf("Burst time of process %d: ", p[i].pid);
        if (scanf("%d", &p[i].burst) != 1 || p[i].burst < 0) {
            fprintf(stderr, "Invalid burst time.\n");
            return 1;
        }

        /* First process starts immediately; subsequent wait for prior finish. */
        p[i].wait = (i == 0) ? 0 : p[i - 1].wait + p[i - 1].burst;
        p[i].turnaround = p[i].wait + p[i].burst;
        total_wait += p[i].wait;
        total_turnaround += p[i].turnaround;
    }

    printf("\nPID  Burst  Wait  Turnaround\n");
    for (i = 0; i < n; i++)
        printf("%3d  %5d  %4d  %10d\n",
               p[i].pid, p[i].burst, p[i].wait, p[i].turnaround);

    printf("\nTotal waiting time    : %d\n", total_wait);
    printf("Average waiting time  : %.2f\n", (double)total_wait / n);
    printf("Total turnaround time : %d\n", total_turnaround);
    printf("Average turnaround    : %.2f\n", (double)total_turnaround / n);

    return 0;
}
\end{lstlisting}

% ============================================================================
% EXPERIMENT 5(d)
% ============================================================================
\exptheader{Experiment 5(d)}{Shortest Job First (SJF) CPU Scheduling}{exp:5.d}

This program simulates non-preemptive Shortest Job First scheduling for processes arriving at $t = 0$. It sorts the process array in non-decreasing order of burst duration using \texttt{qsort()} with an overflow-safe comparator (\texttt{by\_burst}) and accumulates elapsed CPU time to report individual and mean waiting and turnaround times.

\begin{itemize}
    \item \textbf{Waiting-Time Optimality:} Provably minimizes average waiting time among non-preemptive algorithms.
    \item \textbf{Burst Prediction:} Practical kernels estimate next CPU burst via exponential averaging since future bursts are unknown.
    \item \textbf{Starvation Risk:} Continuously arriving short jobs can indefinitely postpone long processes.
\end{itemize}

\begin{lstlisting}[title={\texttt{expt\_5\_d.c}}]
#include <stdio.h>
#include <stdlib.h>

#define MAX_PROCESSES 10

struct process {
    int pid;
    int burst;
    int wait;
    int turnaround;
};

/* Overflow-safe qsort comparator ordering processes by burst time ascending. */
static int by_burst(const void *a, const void *b)
{
    const struct process *x = a;
    const struct process *y = b;

    return (x->burst > y->burst) - (x->burst < y->burst);
}

int main(void)
{
    struct process p[MAX_PROCESSES];
    int n, i, elapsed = 0;
    int total_wait = 0;
    int total_turnaround = 0;

    printf("SJF scheduling\nEnter number of processes: ");
    if (scanf("%d", &n) != 1 || n < 1 || n > MAX_PROCESSES) {
        fprintf(stderr, "Enter a number between 1 and %d.\n", MAX_PROCESSES);
        return 1;
    }

    for (i = 0; i < n; i++) {
        p[i].pid = i + 1;
        printf("Burst time of process %d: ", p[i].pid);
        if (scanf("%d", &p[i].burst) != 1 || p[i].burst < 0) {
            fprintf(stderr, "Invalid burst time.\n");
            return 1;
        }
    }

    /* Sort by shortest burst first to minimize mean waiting time. */
    qsort(p, n, sizeof p[0], by_burst);

    /* Accumulate running clock to compute WT and TAT for each sorted job. */
    for (i = 0; i < n; i++) {
        p[i].wait = elapsed;
        p[i].turnaround = elapsed + p[i].burst;
        elapsed += p[i].burst;
        total_wait += p[i].wait;
        total_turnaround += p[i].turnaround;
    }

    printf("\nPID  Burst  Wait  Turnaround\n");
    for (i = 0; i < n; i++)
        printf("%3d  %5d  %4d  %10d\n",
               p[i].pid, p[i].burst, p[i].wait, p[i].turnaround);

    printf("\nTotal waiting time    : %d\n", total_wait);
    printf("Average waiting time  : %.2f\n", (double)total_wait / n);
    printf("Total turnaround time : %d\n", total_turnaround);
    printf("Average turnaround    : %.2f\n", (double)total_turnaround / n);

    return 0;
}
\end{lstlisting}

% ============================================================================
% EXPERIMENT 6
% ============================================================================
\exptheader{Experiment 6}{Producer--Consumer Problem Using Semaphores}{exp:6}

This program simulates classical bounded-buffer synchronization (capacity $= 3$) between a producer and a consumer through a menu-driven loop. It models a binary semaphore (\texttt{mutex}) for critical-section mutual exclusion alongside counting semaphores (\texttt{empty} and \texttt{full}) manipulated by simulated \texttt{wait\_sem()} ($P$) and \texttt{signal\_sem()} ($V$) operations.

\begin{itemize}
    \item \textbf{Bounded-Buffer Synchronization:} Coordinates asynchronous production and consumption over fixed capacity.
    \item \textbf{Semaphore Primitives:} Binary \texttt{mutex} enforces mutual exclusion; counting \texttt{empty}/\texttt{full} track slot availability.
    \item \textbf{Race Condition Prevention:} Guards buffer state transitions against overflow and underflow.
\end{itemize}

\begin{lstlisting}[title={\texttt{expt\_6.c}}]
#include <stdio.h>

#define BUFFER_SIZE 3

/* Synchronization primitives and shared buffer item counter. */
int mutex = 1;              /* binary semaphore for critical section access */
int full = 0;               /* counting semaphore: occupied buffer slots */
int empty = BUFFER_SIZE;    /* counting semaphore: available buffer slots */
int item = 0;               /* simulated item sequence number */

/* Simulated semaphore P (wait) and V (signal) primitive operations. */
static int wait_sem(int s)   { return s - 1; }
static int signal_sem(int s) { return s + 1; }

/* Enter critical section, reserve an empty slot, produce, and signal full. */
static void produce(void)
{
    mutex = wait_sem(mutex);
    empty = wait_sem(empty);
    full = signal_sem(full);

    printf("Producer produced item %d\n", ++item);

    mutex = signal_sem(mutex);
}

/* Enter critical section, claim a full slot, consume, and signal empty. */
static void consume(void)
{
    mutex = wait_sem(mutex);
    full = wait_sem(full);
    empty = signal_sem(empty);

    printf("Consumer consumed item %d\n", item--);

    mutex = signal_sem(mutex);
}

int main(void)
{
    int choice;

    /* Interactive dispatcher checking buffer preconditions before invocation. */
    while (1) {
        printf("\n1. Produce\n2. Consume\n3. Exit\nChoice: ");
        if (scanf("%d", &choice) != 1)
            break;

        switch (choice) {
        case 1:
            if (empty > 0)
                produce();
            else
                printf("Buffer is full.\n");
            break;
        case 2:
            if (full > 0)
                consume();
            else
                printf("Buffer is empty.\n");
            break;
        case 3:
            return 0;
        default:
            printf("Invalid choice.\n");
        }
    }

    return 0;
}
\end{lstlisting}

% ============================================================================
% EXPERIMENT 8
% ============================================================================
\exptheader{Experiment 8}{Deadlock Avoidance Using Banker's Algorithm}{exp:8}

This program implements Dijkstra's Banker's Safety Algorithm for $n$ processes and $r$ resource types to determine whether the current allocation state is safe. After reading the \texttt{Max}, \texttt{Allocation}, and \texttt{Available} structures, it derives $\text{Need} = \text{Max} - \text{Allocation}$ and iteratively searches for an unfinished process satisfying $\text{Need}_i \le \text{Work}$, reclaiming its resources until either a complete safe sequence is found or an unsafe state is detected.

\begin{itemize}
    \item \textbf{Deadlock Avoidance:} Dynamically checks resource-allocation states rather than statically restricting Coffman conditions.
    \item \textbf{Safe State Invariant:} Guarantees at least one execution sequence allowing all processes to acquire needed resources and terminate.
    \item \textbf{Resource Matrices:} Tracks system state via \texttt{Max}, \texttt{Allocation}, \texttt{Need}, and \texttt{Available} structures.
\end{itemize}

\begin{lstlisting}[title={\texttt{expt\_8.c}}]
#include <stdio.h>
#include <stdlib.h>

#define MAX 100

int n, r;
int max_need[MAX][MAX];
int alloc[MAX][MAX];
int need[MAX][MAX];
int avail[MAX];

/* Read a validated integer from stdin, terminating on malformed input. */
static int read_int(const char *prompt)
{
    int value;

    printf("%s", prompt);
    if (scanf("%d", &value) != 1) {
        fprintf(stderr, "Invalid input.\n");
        exit(EXIT_FAILURE);
    }
    return value;
}

/* Populate an n x r resource matrix (Max or Allocation) from stdin. */
static void read_matrix(int matrix[MAX][MAX], const char *name)
{
    int i, j;

    printf("Enter the %s matrix (%d x %d):\n", name, n, r);
    for (i = 0; i < n; i++) {
        for (j = 0; j < r; j++) {
            if (scanf("%d", &matrix[i][j]) != 1) {
                fprintf(stderr, "Invalid input.\n");
                exit(EXIT_FAILURE);
            }
        }
    }
}

/* Gather system dimensions, Max, Allocation, and Available vectors. */
static void input(void)
{
    int j;

    n = read_int("Enter number of processes: ");
    r = read_int("Enter number of resource types: ");
    if (n < 1 || n > MAX || r < 1 || r > MAX) {
        fprintf(stderr, "Values must be between 1 and %d.\n", MAX);
        exit(EXIT_FAILURE);
    }

    read_matrix(max_need, "Max");
    read_matrix(alloc, "Allocation");

    printf("Enter Available resources (%d values):\n", r);
    for (j = 0; j < r; j++) {
        if (scanf("%d", &avail[j]) != 1) {
            fprintf(stderr, "Invalid input.\n");
            exit(EXIT_FAILURE);
        }
    }
}

/* Display formatted Allocation, Max, Need, and Available resource tables. */
static void show(void)
{
    int i, j;

    printf("\n%-8s %-12s %-12s %-12s %s\n",
           "Process", "Allocation", "Max", "Need", "Available");

    for (i = 0; i < n; i++) {
        printf("P%-7d", i);

        for (j = 0; j < r; j++)
            printf("%3d", alloc[i][j]);
        printf("   ");

        for (j = 0; j < r; j++)
            printf("%3d", max_need[i][j]);
        printf("   ");

        for (j = 0; j < r; j++)
            printf("%3d", need[i][j]);
        printf("   ");

        if (i == 0) {
            for (j = 0; j < r; j++)
                printf("%3d", avail[j]);
        }
        printf("\n");
    }
}

/* Execute Dijkstra's Safety Algorithm to verify existence of a safe sequence. */
static void check_safe(void)
{
    int finish[MAX] = {0};
    int work[MAX];
    int safe[MAX];
    int count = 0;
    int i, j;

    /* Step 1: Compute remaining Need matrix = Max - Allocation. */
    for (i = 0; i < n; i++)
        for (j = 0; j < r; j++)
            need[i][j] = max_need[i][j] - alloc[i][j];

    /* Step 2: Initialize Work vector to currently Available resources. */
    for (j = 0; j < r; j++)
        work[j] = avail[j];

    /* Step 3: Greedily find unfinished processes whose Need <= Work. */
    while (count < n) {
        int found = 0;

        for (i = 0; i < n; i++) {
            if (finish[i])
                continue;

            for (j = 0; j < r && need[i][j] <= work[j]; j++)
                ;

            if (j == r) {               /* all needs can be met */
                /* Simulate process completion and reclaim allocated resources. */
                for (j = 0; j < r; j++)
                    work[j] += alloc[i][j];
                finish[i] = 1;
                safe[count++] = i;
                found = 1;
            }
        }

        if (!found)
            break;                      /* no process can finish */
    }

    /* Step 4: Safe iff every process could be scheduled to completion. */
    if (count == n) {
        printf("\nSystem is in a safe state.\nSafe sequence: ");
        for (i = 0; i < n; i++)
            printf("P%d%s", safe[i], i == n - 1 ? "\n" : " -> ");
    } else {
        printf("\nSystem is in an unsafe state (deadlock possible).\n");
    }
}

int main(void)
{
    printf("Banker's Algorithm\n");
    input();
    check_safe();   /* computes need[][] */
    show();

    return 0;
}
\end{lstlisting}

\end{document}
